\begin{corollary}[Exact transversal number]\label{cor:exacttransversal}
For every $h\ge1$, with $m=2^h$, the host satisfies
\[
 \tau_{H_4}(A_h)=m<m^2=D_{H_4}(A_h).
\]
In particular, the ratio of unrestricted editing distance to the
minimum number of entries hitting all original copies is exactly $m$.
\end{corollary}
\begin{proof}
Let
\[
 \Delta_h=\{(R_h(p),C_h(p)):0\le p<m\}.
\]
Both terminal families contain the same-leaf diagonal
$(R_h(p),C_h(p))$ as their third-row, third-column entry.
The leaf blocks are singletons, so this is one fixed physical cell
for every realization of the corresponding terminal tuple.
Lemma~\ref{lem:classification} therefore gives
$\tau_{H_4}(A_h)\le|\Delta_h|=m$.

For the reverse inequality, use offsets numbered from zero. For
each $p<m$, put $t_p=W^-_{h,p\bmod2}$ and choose a type-I copy
$Q_p$ whose two marker rows, two marker columns and dummy row
all have offset $p$ in their respective blocks. Use the singleton
leaf row and column with label $p$, and use offset $p\bmod2$
in the parent column block $C^-_{h-1}(\lfloor p/2\rfloor)$.
These choices are valid: the marker and dummy blocks have size $m$,
and each parent block has size two. They realize an ordered copy
by the type-I description and Lemma~\ref{lem:quotient}.

The copies $Q_p$ are pairwise disjoint on each axis. Marker
positions of different modes or groups lie in different blocks;
within any one such block, different labels $p$ give different
offsets. The dummy row offsets and leaf labels are also distinct.
Two selected parent columns could agree only if both
$\lfloor p/2\rfloor$ and $p\bmod2$ agree, which determines $p$.
The marker, leaf, dummy and parent categories used on a given axis
are mutually distinct. Thus the $m$ copies are, in particular,
pairwise cell-disjoint. Every set hitting all original copies must
meet each $Q_p$ in a different cell, and hence has size at least $m$.
Combining the bounds proves the transversal identity; the editing
identity is Theorem~\ref{thm:main}, and $m\ge2$ gives strictness.
\end{proof}

The strict gap between hitting original copies and producing a free
matrix sharpens the distinction already developed in \cite{OAI}.
A separate Lean supplement formalizes the exact transversal identity,
using terminal exhaustion from the primary classification theorem.
